\chapter{Local Definitions and Lexical Scope}\label{ch:local}

From this point on we use \textbf{Intermediate Student}, which adds the
\code{local} special form. \code{local} lets us make definitions visible only
inside one expression.

\section{The \code{local} form}
\begin{lstlisting}[style=racket]
(local [(define x1 e1)
        ...
        (define xn en)]
  body)
\end{lstlisting}
The definitions (constants or functions) are visible in \code{body} and in each
other's definitions (later definitions may use earlier ones), but not outside
the \code{local} expression.

\section{Uses of local}

\subsection{1. Clarity: naming intermediate values}
Heron's formula for the area of a triangle with sides $a,b,c$:
$A = \sqrt{s(s-a)(s-b)(s-c)}$ where $s = (a+b+c)/2$.
\begin{lstlisting}[style=racket]
;; (heron a b c) produces the area of a triangle with side lengths a, b, c.
;; heron: Num Num Num -> Num
;;   requires: a, b, c satisfy the triangle inequality
(check-expect (heron 3 4 5) 6)
(define (heron a b c)
  (local [(define s (/ (+ a b c) 2))]
    (sqrt (* s (- s a) (- s b) (- s c)))))
\end{lstlisting}

\subsection{2. Efficiency: avoiding repeated computation}
\begin{lstlisting}[style=racket]
;; list-max: (ne-listof Num) -> Num
(define (list-max lon)
  (cond [(empty? (rest lon)) (first lon)]
        [else (local [(define max-rest (list-max (rest lon)))]
                (cond [(> (first lon) max-rest) (first lon)]
                      [else max-rest]))]))
\end{lstlisting}
Without \code{local}, writing \code{(list-max (rest lon))} in both the
question and the answer makes two recursive calls per level, i.e.\
$\Theta(2^n)$ time. With \code{local} it is evaluated once: $\Theta(n)$.

\subsection{3. Encapsulation: hiding helper functions}
Helpers that are only meaningful to one function can be hidden inside it,
for example the accumulator helper of reverse:
\begin{lstlisting}[style=racket]
;; (my-reverse lst) produces lst in reverse order.
;; my-reverse: (listof Any) -> (listof Any)
(check-expect (my-reverse '(1 2 3)) '(3 2 1))
(define (my-reverse lst)
  (local [;; (rev/acc l acc) produces l reversed, followed by acc
          ;; rev/acc: (listof Any) (listof Any) -> (listof Any)
          (define (rev/acc l acc)
            (cond [(empty? l) acc]
                  [else (rev/acc (rest l) (cons (first l) acc))]))]
    (rev/acc lst empty)))
\end{lstlisting}
Locally defined functions need a purpose and contract, but no separate
examples or tests (they cannot be called from outside, so they are tested
through the enclosing function).

\subsection{4. Using enclosing parameters}
A local function can refer to parameters of the enclosing function, so they
need not be passed along explicitly:
\begin{lstlisting}[style=racket]
;; (count-above lon t) produces how many elements of lon exceed t.
;; count-above: (listof Num) Num -> Nat
(define (count-above lon t)
  (local [;; count/l: (listof Num) -> Nat
          (define (count/l l)
            (cond [(empty? l) 0]
                  [(> (first l) t) (add1 (count/l (rest l)))]
                  [else (count/l (rest l))]))]
    (count/l lon)))
\end{lstlisting}
This becomes essential with \code{lambda} in the next chapter.

\section{Semantics of local}
The substitution model is extended with a rule that \textbf{renames} each
locally defined name to a fresh, unique name and \textbf{lifts} the
definitions to the top level.

\begin{important}[Local rule]
\code{(local [(define x1 e1) ... (define xn en)] body)} $\Rightarrow$
\code{(define x1' e1') ... (define xn' en') body'}, where each \code{xi'} is
a fresh identifier and \code{e1'}, \ldots, \code{body'} are the original
expressions with every \code{xi} replaced by \code{xi'}. The lifted
definitions are then evaluated in order like top-level definitions.
\end{important}

\begin{example}[Trace of local]
\begin{lstlisting}[style=racket]
(heron 3 4 5)
=> (local [(define s (/ (+ 3 4 5) 2))]
     (sqrt (* s (- s 3) (- s 4) (- s 5))))
=> (define s_0 (/ (+ 3 4 5) 2))
   (sqrt (* s_0 (- s_0 3) (- s_0 4) (- s_0 5)))
=> (define s_0 (/ 12 2))
   (sqrt (* s_0 (- s_0 3) (- s_0 4) (- s_0 5)))
=> (define s_0 6)
   (sqrt (* s_0 (- s_0 3) (- s_0 4) (- s_0 5)))
=> (define s_0 6)
   (sqrt (* 6 (- s_0 3) (- s_0 4) (- s_0 5)))
=> ...
=> (define s_0 6)
   (sqrt (* 6 3 2 1))
=> (define s_0 6)
   (sqrt 36)
=> (define s_0 6)
   6
\end{lstlisting}
Every call to \code{heron} creates a new fresh name (\code{s\_0},
\code{s\_1}, \ldots), so values from different calls never collide. In
traces, previously evaluated lifted definitions can be shown once and then
omitted.
\end{example}

\section{Lexical scope}

\begin{definition}[Binding occurrence, scope]
A \textbf{binding occurrence} of an identifier is where it is introduced
(a \code{define}, a function parameter, a \code{local} definition). The
\textbf{scope} of a binding is the region of the program where uses of the
name refer to that binding.
\end{definition}

Racket uses \textbf{lexical (static) scope}: the binding a name refers to is
determined by the program text. To find it, start at the use and move outward
through enclosing \code{local}s and function headers until a binding for that
name is found; if none, it is a top-level (global) binding.

\begin{example}[Shadowing]
\begin{lstlisting}[style=racket]
(define x 1)
(define (f x)
  (local [(define x 10)
          (define (g y) (+ x y))]
    (g x)))
(f 5)   ; => 20
\end{lstlisting}
Inside \code{g}, \code{x} refers to the innermost binding, the local
\code{x} = 10, which \textbf{shadows} the parameter \code{x} = 5 and the global
\code{x} = 1. The argument in \code{(g x)} is also the local \code{x}. So the
result is $10 + 10 = 20$.
\end{example}

It helps to draw \textbf{contour diagrams}: a box around each scope (top
level, each function body, each \code{local}); a use refers to the nearest
enclosing box that binds the name.

\section{Style}
\begin{itemize}
  \item Use \code{local} to name values that are used more than once or that have a meaningful name.
  \item Do not over-nest: if a helper is useful on its own, keep it at top level and test it.
  \item Avoid reusing the same name at different levels (shadowing) unless deliberate; it confuses readers.
\end{itemize}

\section{Summary}
\begin{itemize}
  \item \code{local} provides definitions visible in a single expression.
  \item Uses: clarity, efficiency (avoid recomputation), encapsulation, and access to enclosing parameters.
  \item Semantics: rename to fresh names and lift to the top level.
  \item Racket is lexically scoped; inner bindings shadow outer ones.
\end{itemize}
